% !PRINTED 596
\weberpage{596}
%<<CONT
is. The number $\boldsymbol{s}(u)^2$ generates the field $\mathfrak{T}_{\mathfrak{m}}$. Moreover,
by § 122, for every number $\varrho$ in $\mathfrak{L}$:

\weq{11}{\varrho \equiv \varrho^{p} \quad (\mathrm{mod}\ \mathfrak{P}),}

if $\mathfrak{P}$ is a prime ideal in $\mathfrak{L}$ that divides $\pi$, and consequently,
for every integer $\tau$ in $\mathfrak{T}_{\mathfrak{m}}$, if $\mathfrak{P}'$ is a prime ideal in $\mathfrak{T}_{\mathfrak{m}}$:

\weq{12}{\tau^{p} \equiv \tau \quad (\mathrm{mod}\ \mathfrak{P}'),}

if the conditions (10) and (1) are satisfied.

From this it follows:

8. \gesperrt{In order that a prime $\pi$ existing in $\Omega$ should split in the field $\mathfrak{T}_{\mathfrak{m}}$ into prime ideals of the first degree, it is necessary and sufficient that}
\[
\begin{aligned}
\pi &\equiv 1 \quad (\mathrm{mod}\ 2), \\
\pi &\equiv \pm 1 \quad (\mathrm{mod}\ \mathfrak{m}).
\end{aligned}
\]

A certain finite number of primes (those which divide the
discriminants of the equations that occur) is here left out of
account.

\begin{center}
\textbf{\S\,161.\qquad \gesperrt{Number Groups and Ideal Groups.}}
\end{center}

As in § 98 and § 106 I denote by $O$ the totality of the
integral and fractional numbers of the quadratic field $\Omega$,
after the exclusion of all numbers that are not coprime to an
arbitrarily assumed $S$. $S$ may be a rational number, a number
in $\Omega$, or an ideal as well. I will call it the \gesperrt{excludent}.
In the three memoirs \glqq über Zahlgruppen in algebraischen
 Körpern\grqq{}\,\textsuperscript{1)} I have gathered together, under the name \gesperrt{number groups},
certain systems of numbers out of $O$ which were defined by
the requirement that the \gesperrt{product} and the \gesperrt{quotient} of any two
numbers of such a group were contained in the same group.
Such a number group can never contain the number zero,
but certainly contains the number 1\,\textsuperscript{2)}.

In the same sense in which we speak of number groups we
may also speak of \gesperrt{ideal groups}, which likewise may have their
excludent. Here fractional ideals naturally occur as well, and
in dealing with them one makes use, not without advantage, of the mode of representation
 by means of the \gesperrt{functionals} (Bd. II, § 169).

\begingroup
\footnotesize
\noindent\rule{4cm}{0.4pt}\par
\textsuperscript{1)}\, Mathematische Annalen Bd. 48, 49, 50.

\textsuperscript{2)}\, \gesperrt{Fueter} calls these groups \glqq \gesperrt{Zahlstrahlen}\grqq{}. Cf. \gesperrt{Fueter}, Die
Theorie der Zahlstrahlen I, II; \gesperrt{Crelle}, Bd. 130, 132.
\endgroup
